% Generated draft: requires mathematical and source review.
\documentclass[11pt]{article}
\usepackage[margin=25mm]{geometry}
\usepackage{fontspec}
\setmainfont{Latin Modern Roman}
\usepackage{amsmath,amssymb,mathtools,unicode-math}
\setmathfont{Latin Modern Math}
\usepackage{enumitem,needspace,xurl,hyperref}
\setlength{\parindent}{0pt}
\setlength{\parskip}{5pt}
\setlength{\emergencystretch}{3em}
\newcommand{\N}{\mathbb{N}}
\newcommand{\Z}{\mathbb{Z}}
\newcommand{\Q}{\mathbb{Q}}
\newcommand{\R}{\mathbb{R}}
\newcommand{\C}{\mathbb{C}}
\newcommand{\F}{\mathbb{F}}
\newcommand{\A}{\mathbb{A}}
\newcommand{\PP}{\mathbb{P}}
\newcommand{\E}{\mathbb{E}}
\newcommand{\eps}{\varepsilon}
\newcommand{\dd}{\,\mathrm{d}}
\begin{document}
\section*{UnsolvedMath ID: 30000564}
\textbf{Source problem code:} catalogue entry
\subsection{Definitions and mathematical statement}
Let \(M\) be a Haken three-manifold in the source's intended convention, and let \(\Gamma=\pi_1(M)\). In the usual compact orientable setting, Haken refers to irreducibility together with a suitable properly embedded two-sided incompressible surface; the precise boundary and sufficiently-large conventions are not stated in the record. The subgroup question below is independent of choosing a particular presentation of \(\Gamma\).

A subgroup \(H\leq\Gamma\) is separable in \(\Gamma\) when
\[
 \forall\gamma\in\Gamma\setminus H\ \exists Q\text{ finite}\ \exists q:\Gamma\to Q
 \text{ a homomorphism},\qquad q(\gamma)\notin q(H).
\]
Equivalently, it is closed in the profinite topology on \(\Gamma\), or
\[
 H=\bigcap_{\substack{H\leq K\leq\Gamma\\[\Gamma:K]<\infty}}K.
\]
The finite-index supergroups \(K\) in this intersection are not required to be normal. The finite quotient or supergroup may depend on the element \(\gamma\) being excluded.

The question is to characterize those subgroups \(H\leq\pi_1(M)\) which have this property, across the intended Haken manifolds. The extracted text does not restrict \(H\) to finitely generated subgroups, surface subgroups, or subgroups associated to particular geometric pieces. Such restrictions would change the classification and are not inserted silently. Nor does ``which'' explicitly demand an algorithm from a chosen finite input, as opposed to structural necessary and sufficient conditions.

Separability is a property of the embedding \(H\leq\Gamma\), not merely residual finiteness of \(H\) as an abstract group. Residual finiteness of \(\Gamma\) tests the case \(H=\{1\}\); it is not a replacement for separation of arbitrary subgroups. No assertion that every subgroup, or every finitely generated subgroup, is separable is built into the question.

The literature triage checked 2026-08-21 reports no authoritative resolution of the extracted claim located in its targeted search. The exact-statement label says the sentence was retained byte-for-byte. These remain dated source metadata, not an independently verified current classification or a denial of results for narrower manifold and subgroup classes.
\subsection{Short English statement}
Characterize subgroups of Haken three-manifold groups that are closed in the ambient profinite topology. Each outside element must be excluded from the subgroup's image in some finite quotient. The source does not restrict to finitely generated subgroups or request a particular algorithm.
\subsection{Sources}
\begin{itemize}
\item Local UnsolvedMath registry, permanent ID 30000564, source code catalogue entry, ``Separable Subgroups of Haken Three-Manifold Groups,'' including the unchanged original question and exact-statement verification.
\item \emph{Four-dimensional Manifolds (2007)}, Oberwolfach Reports, pp.\ 2141--2194, DOI 10.4171/owr/2006/35: \url{https://doi.org/10.4171/owr/2006/35}. This is also the sole supplied literature reference. The citation uses 2007, while the DOI and proposed-year metadata use 2006.
\item Literature triage checked 2026-08-21; record updated 2026-09-16. No external source has been visited or independently checked here.
\end{itemize}
\end{document}
